\section{Diseño del Sistema y Arquitectura}
\label{sec:arquitectura}

\subsection{Modelo C4}

La arquitectura se documenta en tres niveles del modelo C4 (\texttt{docs/arquitectura/README.md}), cuyos
diagramas se muestran en las Figuras~\ref{fig:c4-contexto}, \ref{fig:c4-contenedores} y
\ref{fig:c4-componentes} --- se mantienen en inglés, como el resto de esa documentación, para que sean
legibles/indexables de forma independiente del texto circundante en español:

\begin{itemize}[nosep]
    \item \textbf{Nivel 1 (Contexto, Figura~\ref{fig:c4-contexto}):} el sistema se relaciona con cuatro tipos de actores (Estudiante,
    Docente/Jurado, Coordinador, Administrador) y un sistema externo (API de universidades de Hipo Labs,
    consumida vía \texttt{ExternalApiServiceImpl} con caché Redis de 10 minutos).
    \item \textbf{Nivel 2 (Contenedores, Figura~\ref{fig:c4-contenedores}):} cuatro contenedores desplegables --- Frontend Angular (SPA),
    Backend Spring Boot (API REST), PostgreSQL 15 (persistencia), Redis 7 (caché), coordinados por nginx
    como proxy inverso.
    \item \textbf{Nivel 3 (Componentes, Figura~\ref{fig:c4-componentes}):} 31 controladores REST, organizados por dominio (Auth,
    Solicitudes, Jurados, Cronograma, Evaluaciones, Actas, Notificaciones, Reportes, Salas, Usuarios,
    Anteproyectos, y otros de soporte), cada uno delegando a una capa de servicio y, para la lógica
    académica compleja, a procedimientos almacenados PL/pgSQL.
\end{itemize}

\begin{figure}[H]
\centering
\includegraphics[width=0.85\textwidth]{figuras/fig-c4-contexto}
\caption{C4 Nivel 1 --- Diagrama de contexto del sistema.}
\label{fig:c4-contexto}
\end{figure}

\begin{figure}[H]
\centering
\includegraphics[width=0.85\textwidth]{figuras/fig-c4-contenedores}
\caption{C4 Nivel 2 --- Diagrama de contenedores.}
\label{fig:c4-contenedores}
\end{figure}

\begin{figure}[H]
\centering
\includegraphics[width=0.95\textwidth]{figuras/fig-c4-componentes}
\caption{C4 Nivel 3 --- Diagrama de componentes del backend.}
\label{fig:c4-componentes}
\end{figure}

\subsection{Resumen de los siete Registros de Decisión Arquitectónica (ADR)}

\begin{table}[H]
\centering
\small
\begin{tabular}{p{1.6cm}p{3.2cm}p{7cm}}
\toprule
\textbf{ADR} & \textbf{Decisión} & \textbf{Justificación resumida} \\
\midrule
ADR-001 & Arquitectura general de 3 capas & Separación Angular/Spring Boot/PostgreSQL para independencia de despliegue y escalado \\
ADR-002 & Autenticación JWT + cookies HTTP-Only & Stateless para escalar horizontalmente sin sesión compartida; cookie HTTP-Only mitiga XSS sobre el token \\
ADR-003 & Frontend Angular (SPA) & Reactividad y componentización para un dominio con múltiples roles y vistas condicionales \\
ADR-004 & Seguridad alineada a OWASP & JWT de 7 claims, rate limiting, cabeceras HSTS/CSP/X-Frame-Options, BCrypt \\
ADR-005 & Control de acceso por permisos dinámicos (BD), no roles estáticos en código & Reemplaza el RBAC fijo con \texttt{@PreAuthorize} descrito originalmente en ADR-004 por una verificación contra \texttt{permisos}/\texttt{rol\_permisos}, para que cambiar qué rol puede qué requiera solo datos, no redespliegue \\
ADR-006 & Estrategia híbrida JPA + procedimientos almacenados (separación CRUD/SP) & CRUD simple vía Spring Data JPA; lógica académica compleja (notas ponderadas, reportes, firmas) vía PL/pgSQL para evitar SQL dinámico. Ampliada en la Fase 8 con la decisión explícita de mantener la convención Flyway en vez de espejar en \texttt{db/procs/} \\
ADR-007 & Despliegue con Docker Compose, imágenes ancladas por digest SHA-256 & Reproducibilidad determinista. Ampliada en la Fase 8 con la decisión de Railway como proveedor de producción (Sección~\ref{sec:implementacion}) \\
\bottomrule
\end{tabular}
\caption{Summary of the repository's 7 real ADRs (\texttt{docs/adr/}).}
\label{tab:adrs}
\end{table}

\textbf{Corrección de integridad relevante:} ADR-007 afirmaba, en su redacción original, que las
imágenes base ya estaban ancladas por digest SHA-256 en \texttt{docker-compose.yml}. Verificado contra el
archivo real durante la Fase 8 de este proyecto: esa afirmación era falsa (las imágenes usaban solo tags
mutables). Se implementó el anclaje real en esa misma revisión, con los tres digests verificados vía
\texttt{docker pull}, en vez de simplemente corregir el texto de la ADR sin corregir el código que
describía. (Nota: ADR-006 y ADR-007 se renumeraron el 2026-09-01 --- antes eran ADR-003 y ADR-006
respectivamente --- para que coincidan con los temas específicos que exige la guía de la Entrega Final
en esos dos números; el contenido no cambió, solo el orden.)

\subsection{Modelo de calidad (ISO/IEC 25010)}

La Tabla~\ref{tab:iso25010} mapea las ocho características del modelo de calidad de producto de software
ISO/IEC 25010:2011 \cite{iso25010} contra la evidencia empírica real de este proyecto, en vez de afirmar
cumplimiento genérico sin sustento. Las cifras remiten a la Sección~\ref{sec:evaluacion} (evaluación
empírica) y a la Sección~\ref{sec:seguridad-corregida} (correcciones de seguridad).

\begin{table}[H]
\centering
\small
\begin{tabular}{p{2.6cm}p{9.5cm}}
\toprule
\textbf{Característica} & \textbf{Evidencia real en este proyecto} \\
\midrule
Adecuación funcional & 36/50 requisitos Must verificados (72\,\%); 90 requisitos totales (64
funcionales + 26 no funcionales); matriz de trazabilidad de 90 filas validada contra el código real por
\texttt{validate-traceability.sh} (Sección~\ref{sec:requisitos}). \\
Eficiencia de desempeño & p95 de \texttt{http\_req\_duration} entre 6.17 y 8.53\,ms en caché caliente
(5 corridas k6 independientes); prueba de Mann-Whitney U entre caché fría/caliente con
$p = 3.02\times10^{-11}$ (Sección~\ref{sec:evaluacion}). \\
Compatibilidad & API REST versionada (\texttt{/api/v1}) documentada con OpenAPI 3.0/Swagger UI; JSON como
formato de intercambio estándar en los 31 controladores. \\
Usabilidad & Instrumento SUS \cite{brooke1996} aplicado a 15 personas; \textbf{solo 4 con fecha de
aplicación verificable} ($n=4$, media 48.75/100) --- las 11 restantes quedan declaradas pendientes de
confirmar en vez de contarse como cerradas (Sección~\ref{sec:evaluacion}). \\
Fiabilidad & 559 pruebas automatizadas reales (0 fallos), \texttt{/actuator/health} como verificación de
arranque, migraciones Flyway versionadas y reproducibles end-to-end. \\
Seguridad & Auditoría OWASP ZAP (0 hallazgos High), SpotBugs/find-sec-bugs (0 SQL dinámico), 6 fallos
reales de autorización (IDOR/mass-assignment) encontrados y corregidos por revisión manual dirigida
(Sección~\ref{sec:seguridad-corregida}). \\
Mantenibilidad & Arquitectura en capas (controlador/servicio/repositorio) con 31 controladores
organizados por dominio; 7 ADRs documentando decisiones y alternativas descartadas (Tabla~\ref{tab:adrs}). \\
Portabilidad & Despliegue reproducible con Docker Compose y las tres imágenes base ancladas por digest
SHA-256 (ADR-007); \textbf{sin despliegue público todavía} --- brecha declarada explícitamente
(Sección~\ref{sec:trabajo-futuro}). \\
\bottomrule
\end{tabular}
\caption{Características de calidad ISO/IEC 25010 mapeadas a evidencia real del proyecto.}
\label{tab:iso25010}
\end{table}

\subsection{Modelo de datos}

El esquema relacional (\texttt{presus}, PostgreSQL 15) modela 13 entidades principales (usuarios,
solicitudes, anteproyectos, jurados, cronogramas, evaluaciones, actas, rúbricas, salas, notificaciones,
entre otras) con claves foráneas que garantizan integridad referencial. Los identificadores son
\texttt{BIGINT} autoincrementales (migrados desde UUID en una corrección temprana documentada en
OBS-01), y diez rutinas PL/pgSQL puras (ocho procedimientos/funciones invocados desde Java, más dos
funciones de trigger de auditoría que corren solo a nivel de motor --- catálogo completo en
\texttt{docs/basedatos/CATALOGO-SP.md}) encapsulan la lógica académica que involucra uniones
multi-tabla, versionados vía Flyway en \texttt{backend/src/main/resources/db/migration/}.

El Listado~\ref{lst:sp-firmar-acta} muestra uno de esos procedimientos, \texttt{sp\_firmar\_acta\_digital}
(firma digital multi-actor de actas), y el Listado~\ref{lst:acta-service-firma} el método
\texttt{MinutesServiceImpl.signMinutes()} que lo invoca vía JPA 2.1 \texttt{@Procedure}
(\texttt{MinutesRepository.signMinutesDigital}).

\begin{lstlisting}[language=SQL,style=codigobase,caption={\texttt{sp\_firmar\_acta\_digital}: procedimiento almacenado PL/pgSQL (migracion V10).},label={lst:sp-firmar-acta}]
CREATE OR REPLACE PROCEDURE presus.sp_firmar_acta_digital(
    p_acta_id BIGINT,
    p_rol VARCHAR,
    p_observacion TEXT
) AS $$
BEGIN
    IF UPPER(p_rol) = 'PRESIDENTE' THEN
        UPDATE presus.actas
        SET firmada_presidente = true,
            fecha_firma_presidente = NOW(),
            observaciones_acta = COALESCE(observaciones_acta, '') ||
                E'\n[PRESIDENTE]: ' || COALESCE(p_observacion, 'Firma registrada')
        WHERE id = p_acta_id;
    ELSIF UPPER(p_rol) = 'VOCAL_1' OR UPPER(p_rol) = 'VOCAL1' THEN
        UPDATE presus.actas
        SET firmada_vocal1 = true,
            fecha_firma_vocal1 = NOW(),
            observaciones_acta = COALESCE(observaciones_acta, '') ||
                E'\n[VOCAL_1]: ' || COALESCE(p_observacion, 'Firma registrada')
        WHERE id = p_acta_id;
    -- ... ramas analogas para VOCAL_2 y TUTOR (omitidas por espacio) ...
    ELSE
        RAISE EXCEPTION 'Rol invalido para firma de acta: %', p_rol;
    END IF;

    -- Consolidar el estado general de firmas
    UPDATE presus.actas
    SET firmada = (firmada_presidente AND firmada_vocal1 AND firmada_vocal2 AND firmada_tutor)
    WHERE id = p_acta_id;
END;
$$ LANGUAGE plpgsql;
\end{lstlisting}

\begin{lstlisting}[language=Java,style=codigobase,caption={Invocacion real desde Java (repositorio JPA 2.1 \texttt{@Procedure} y servicio que lo llama).},label={lst:acta-service-firma}]
// MinutesRepository.java
@Procedure(procedureName = "presus.sp_firmar_acta_digital")
void signMinutesDigital(@Param("p_acta_id") Long minutesId,
                        @Param("p_rol") String role,
                        @Param("p_observacion") String observacion);

// MinutesServiceImpl.java
@Override
@Transactional
public Minutes signMinutes(Long minutesId, String role, String observacion) {
    // ... validacion de rol y de que el appUser autenticado es
    // el presidente/vocal/tutor correspondiente (omitido por espacio) ...

    // Persiste la firma via sp_firmar_acta_digital (Fase 3 / Criterio P1) -- el
    // procedimiento es la fuente de verdad de firmada_*/fecha_firma_*/observaciones_minutes;
    // se refresca la entidad para que el resto del metodo vea lo que el SP
    // realmente persistio, en vez de sobreescribirlo con el save() final.
    minutesRepository.signMinutesDigital(minutesId, roleNormalizado, observacion);
    entityManager.refresh(minutes);

    acta.actualizarEstadoFirma();
    // ... notificacion al estudiante y persistencia final (omitido) ...
    return acta;
}
\end{lstlisting}

Un segundo caso, más complejo, es el de un procedimiento que devuelve un conjunto de filas
(\texttt{sp\_calcular\_promedio\_evaluacion}): PostgreSQL rechaza la sintaxis \texttt{CALL} que Hibernate
genera contra una \texttt{FUNCTION}, así que la única forma de que la firma coincida es declararlo como
\texttt{PROCEDURE} con un parámetro \texttt{INOUT refcursor}, mapeado en Java con
\texttt{ParameterMode.REF\_CURSOR} (Listado~\ref{lst:namedstoredprocedure}).

\begin{lstlisting}[language=Java,style=codigobase,caption={Mapeo JPA 2.1 de un procedimiento con retorno de conjunto de filas via \texttt{REF\_CURSOR} (\texttt{Evaluacion.java}).},label={lst:namedstoredprocedure}]
@NamedStoredProcedureQuery(
        name = "Evaluacion.calcularPromedioEvaluacion",
        procedureName = "presus.sp_calcular_promedio_evaluacion",
        resultSetMappings = "PromedioEvaluacionMapping",
        parameters = {
                @StoredProcedureParameter(mode = ParameterMode.IN,
                        name = "p_solicitud_id", type = Long.class),
                @StoredProcedureParameter(mode = ParameterMode.REF_CURSOR,
                        name = "p_resultado", type = void.class)
        }
)
@SqlResultSetMapping(
        name = "PromedioEvaluacionMapping",
        classes = @ConstructorResult(
                targetClass = ec.edu.uteq.presustentaciones.dto.PromedioEvaluacionResult.class,
                columns = {
                        @ColumnResult(name = "solicitud_id", type = Long.class),
                        @ColumnResult(name = "nota_final", type = Double.class),
                        @ColumnResult(name = "estado_resultado", type = String.class)
                }
        )
)
\end{lstlisting}
\newpage
